% GENERATED FILE. Edit canonical lessons, then run npm run generate:books.
% canonical-source-hash: fnv1a64:3c66123907bb4acb
% canonical-lessons: JA-C526-shinpo-suru, JA-C526-hatten-suru, JA-C526-kaizen-suru, JA-C526-kouten-suru, JA-C526-dondon

\chapter{Progress}
\label{ch:ja-progress}
\hlchaptermodality{526}

\begin{chapteropening}
\textbf{By the end of this chapter:} \emph{I can say how things would get better.}

Complete the last lesson of chapter 526: I can say how things would get better.
\end{chapteropening}

\section[\texorpdfstring{\ja{しんぽする} (shinpo suru)}{しんぽする (shinpo suru)}]{\ja{しんぽする} (shinpo suru) — to progress; to advance}

\label{lesson:JA-C526-shinpo-suru}



\begin{quote}
Before the new one: say the Japanese for to compare (to), then the Japanese for sooner or later.
\end{quote}

\subsection*{You'll want to know: \ja{しんぽする}}
\textbf{\ja{しんぽする}} — \emph{shinpo suru} — \textquotedblleft{}to progress; to advance\textquotedblright{}.

\begin{quote}
\textbf{\ja{もし} \ja{ぎじゅつが} \ja{しんぽしたら、くるまは} \ja{はやく} \ja{なるでしょう。}} — \emph{moshi gijutsu ga shinpo shitara, kuruma wa hayaku naru deshō} — \textquotedblleft{}If technology advances, cars will get faster.\textquotedblright{}
\end{quote}

Polite \textbf{\ja{しんぽします}}, \emph{shinposhimasu}; a \ja{する} verb. Past (polite): \textbf{\ja{ぎじゅつが} \ja{しんぽしました。}}, \emph{gijutsu ga shinposhimashita}.

\begin{cousinweb}
\textbf{\ja{しんぽする}} (\emph{shinpo suru}) is the Sino-Japanese (kango) noun \emph{shinpo} with the native verb \emph{suru}, \textquotedblleft{}to do\textquotedblright{}: two Chinese characters read with their on'yomi, readings borrowed from Chinese. They are \emph{shin}, \textquotedblleft{}advance\textquotedblright{} (Middle Chinese \emph{tsinH}) and \emph{po}, \textquotedblleft{}walk\textquotedblright{} (KANJIDIC2; Middle Chinese in Baxter's transcription).
\end{cousinweb}

\begin{grammarlens}[title={What you've built}]
One more word for saying what happens if.
\end{grammarlens}

\subsection*{Your turn}
\begin{itemize}
\raggedright
  \item \emph{Say it:} \emph{shinpo suru}
  \item \emph{Say it:} \emph{shinpo suru}, once more
  \item \emph{Say it:} \emph{moshi gijutsu ga shinpo shitara, kuruma wa hayaku naru deshō} — If technology advances, cars will get faster.
  \item \emph{Recall:} say the Japanese for to compare (to), then the Japanese for sooner or later, then say \emph{shinpo suru} again
\end{itemize}

\begin{tcolorbox}[breakable,colback=teal!4,colframe=teal!35!black,title={Before you move on}]
What does \ja{しんぽする} mean? (\textquotedblleft{}To progress; to advance\textquotedblright{}.) Say it once more.
\end{tcolorbox}
\section[\texorpdfstring{\ja{はってんする} (hatten suru)}{はってんする (hatten suru)}]{\ja{はってんする} (hatten suru) — to develop; to grow}

\label{lesson:JA-C526-hatten-suru}



\begin{quote}
Before the new one: say the Japanese for sooner or later, then the Japanese for to progress.
\end{quote}

\subsection*{You'll want to know: \ja{はってんする}}
\textbf{\ja{はってんする}} — \emph{hatten suru} — \textquotedblleft{}to develop; to grow\textquotedblright{}.

\begin{quote}
\textbf{\ja{もし} \ja{こうじょうが} \ja{きたら、まちは} \ja{はってんするでしょう。}} — \emph{moshi kōjō ga kitara, machi wa hatten suru deshō} — \textquotedblleft{}If a factory comes, the town will probably grow.\textquotedblright{}
\end{quote}

Polite \textbf{\ja{はってんします}}, \emph{hattenshimasu}; a \ja{する} verb. Past (polite): \textbf{\ja{まちが} \ja{はってんしました。}}, \emph{machi ga hattenshimashita}.

\begin{cousinweb}
\textbf{\ja{はってんする}} (\emph{hatten suru}) is the Sino-Japanese (kango) noun \emph{hatten} with the native verb \emph{suru}, \textquotedblleft{}to do\textquotedblright{}: two Chinese characters read with their on'yomi, readings borrowed from Chinese. They are \emph{ha}, \textquotedblleft{}departure\textquotedblright{} (Middle Chinese \emph{pjot}) and \emph{ten}, \textquotedblleft{}unfold\textquotedblright{} (\emph{trjenX}) (KANJIDIC2; Middle Chinese in Baxter's transcription).
\end{cousinweb}

\begin{grammarlens}[title={What you've built}]
One more word for saying what happens if.
\end{grammarlens}

\subsection*{Your turn}
\begin{itemize}
\raggedright
  \item \emph{Say it:} \emph{hatten suru}
  \item \emph{Say it:} \emph{hatten suru}, once more
  \item \emph{Say it:} \emph{moshi kōjō ga kitara, machi wa hatten suru deshō} — If a factory comes, the town will probably grow.
  \item \emph{Recall:} say the Japanese for sooner or later, then the Japanese for to progress, then say \emph{hatten suru} again
\end{itemize}

\begin{tcolorbox}[breakable,colback=teal!4,colframe=teal!35!black,title={Before you move on}]
What does \ja{はってんする} mean? (\textquotedblleft{}To develop; to grow\textquotedblright{}.) Say it once more.
\end{tcolorbox}
\section[\texorpdfstring{\ja{かいぜんする} (kaizen suru)}{かいぜんする (kaizen suru)}]{\ja{かいぜんする} (kaizen suru) — to improve (something)}

\label{lesson:JA-C526-kaizen-suru}



\begin{quote}
Before the new one: say the Japanese for to progress, then the Japanese for to develop.
\end{quote}

\subsection*{You'll want to know: \ja{かいぜんする}}
\textbf{\ja{かいぜんする}} — \emph{kaizen suru} — \textquotedblleft{}to improve (something)\textquotedblright{}.

\begin{quote}
\textbf{\ja{もし} \ja{もんだいを} \ja{かいぜんしたら、きゃくが} \ja{ふえるでしょう。}} — \emph{moshi mondai o kaizenshitara, kyaku ga fueru deshō} — \textquotedblleft{}If we fix the problems, customers will probably increase.\textquotedblright{}
\end{quote}

Polite \textbf{\ja{かいぜんします}}, \emph{kaizenshimasu}; a \ja{する} verb. Past (polite): \textbf{\ja{てんちょうが} \ja{みせを} \ja{かいぜんしました。}}, \emph{tenchō ga mise o kaizenshimashita}.

\begin{cousinweb}
\textbf{\ja{かいぜんする}} (\emph{kaizen suru}) is the Sino-Japanese (kango) noun \emph{kaizen} with the native verb \emph{suru}, \textquotedblleft{}to do\textquotedblright{}: two Chinese characters read with their on'yomi, readings borrowed from Chinese. They are \emph{kai}, \textquotedblleft{}reformation\textquotedblright{} (Middle Chinese \emph{kojX}) and \emph{zen}, \textquotedblleft{}virtuous\textquotedblright{} (\emph{dzyenX}) (KANJIDIC2; Middle Chinese in Baxter's transcription).
\end{cousinweb}

\begin{grammarlens}[title={What you've built}]
One more word for saying what happens if.
\end{grammarlens}

\subsection*{Your turn}
\begin{itemize}
\raggedright
  \item \emph{Say it:} \emph{kaizen suru}
  \item \emph{Say it:} \emph{kaizen suru}, once more
  \item \emph{Say it:} \emph{moshi mondai o kaizenshitara, kyaku ga fueru deshō} — If we fix the problems, customers will probably increase.
  \item \emph{Recall:} say the Japanese for to progress, then the Japanese for to develop, then say \emph{kaizen suru} again
\end{itemize}

\begin{tcolorbox}[breakable,colback=teal!4,colframe=teal!35!black,title={Before you move on}]
What does \ja{かいぜんする} mean? (\textquotedblleft{}To improve (something)\textquotedblright{}.) Say it once more.
\end{tcolorbox}
\section[\texorpdfstring{\ja{こうてんする} (kōten suru)}{こうてんする (kōten suru)}]{\ja{こうてんする} (kōten suru) — to turn for the better}

\label{lesson:JA-C526-kouten-suru}



\begin{quote}
Before the new one: say the Japanese for to develop, then the Japanese for to improve (something).
\end{quote}

\subsection*{You'll want to know: \ja{こうてんする}}
\textbf{\ja{こうてんする}} — \emph{kōten suru} — \textquotedblleft{}to turn for the better\textquotedblright{}.

\begin{quote}
\textbf{\ja{もし} \ja{てんきが} \ja{こうてんしたら、しあいを} \ja{します。}} — \emph{moshi tenki ga kōtenshitara, shiai o shimasu} — \textquotedblleft{}If the weather improves, we will hold the match.\textquotedblright{}
\end{quote}

Polite \textbf{\ja{こうてんします}}, \emph{kōtenshimasu}; a \ja{する} verb. Past (polite): \textbf{\ja{びょうきが} \ja{こうてんしました。}}, \emph{byōki ga kōtenshimashita}.

\begin{cousinweb}
\textbf{\ja{こうてんする}} (\emph{kōten suru}) is the Sino-Japanese (kango) noun \emph{kōten} with the native verb \emph{suru}, \textquotedblleft{}to do\textquotedblright{}: two Chinese characters read with their on'yomi, readings borrowed from Chinese. They are \emph{kō}, \textquotedblleft{}pleasing\textquotedblright{} (Middle Chinese \emph{xawX}) and \emph{ten}, \textquotedblleft{}revolve\textquotedblright{} (\emph{trjwenX}) (KANJIDIC2; Middle Chinese in Baxter's transcription).
\end{cousinweb}

\begin{grammarlens}[title={What you've built}]
One more word for saying what happens if.
\end{grammarlens}

\subsection*{Your turn}
\begin{itemize}
\raggedright
  \item \emph{Say it:} \emph{kōten suru}
  \item \emph{Say it:} \emph{kōten suru}, once more
  \item \emph{Say it:} \emph{moshi tenki ga kōtenshitara, shiai o shimasu} — If the weather improves, we will hold the match.
  \item \emph{Recall:} say the Japanese for to develop, then the Japanese for to improve (something), then say \emph{kōten suru} again
\end{itemize}

\begin{tcolorbox}[breakable,colback=teal!4,colframe=teal!35!black,title={Before you move on}]
What does \ja{こうてんする} mean? (\textquotedblleft{}To turn for the better\textquotedblright{}.) Say it once more.
\end{tcolorbox}
\section[\texorpdfstring{\ja{どんどん} (dondon)}{どんどん (dondon)}]{\ja{どんどん} (dondon) — rapidly; steadily (and fast)}

\label{lesson:JA-C526-dondon}



\begin{quote}
Before the new one: say the Japanese for to improve (something), then the Japanese for to turn for the better.
\end{quote}

\subsection*{You'll want to know: \ja{どんどん}}
\textbf{\ja{どんどん}} — \emph{dondon} — \textquotedblleft{}rapidly; steadily (and fast)\textquotedblright{}.

\begin{quote}
\textbf{\ja{もし} \ja{あめが} \ja{ふったら、かわの} \ja{みずが} \ja{どんどん} \ja{ふえるでしょう。}} — \emph{moshi ame ga futtara, kawa no mizu ga dondon fueru deshō} — \textquotedblleft{}If it rains, the river will rise fast.\textquotedblright{}
\end{quote}

\begin{cousinweb}
\textbf{\ja{どんどん}} (\emph{dondon}) is a native mimetic word, imitating a sound rather than borrowed from Chinese: JMdict's first sense is the noise of drumming, and its later senses are \textquotedblleft{}rapidly\textquotedblright{} and \textquotedblleft{}steadily\textquotedblright{}; no dictionary we use records an older form, so its history is not settled (JMdict).
\end{cousinweb}

\begin{grammarlens}[title={What you've built}]
One more word for saying what happens if.
\end{grammarlens}

\subsection*{Your turn}
\begin{itemize}
\raggedright
  \item \emph{Say it:} \emph{dondon}
  \item \emph{Say it:} \emph{dondon}, once more
  \item \emph{Say it:} \emph{moshi ame ga futtara, kawa no mizu ga dondon fueru deshō} — If it rains, the river will rise fast.
  \item \emph{Recall:} say the Japanese for to improve (something), then the Japanese for to turn for the better, then say \emph{dondon} again
\end{itemize}

\begin{tcolorbox}[breakable,colback=teal!4,colframe=teal!35!black,title={Before you move on}]
What does \ja{どんどん} mean? (\textquotedblleft{}Rapidly; steadily (and fast)\textquotedblright{}.) Say it once more.
\end{tcolorbox}
